% Propietario conservado: /Users/ruben/Documents/ChatGPT/jueces y controles/REVISION_ARGUMENTAL_APERTURAS_CIERRES_20260915/IV/ES/sections/acoplamiento_recuperacion.tex
\section{Arithmetic--geometric coupling and bilateral recovery}
\label{iv:sec:acoplamiento-recuperacion}

The dodecaphase register of section~\ref{iv:sec:k-registro} admits
arithmetic, incidence and operator readings. We call
\emph{holographic constant of arithmetic-geometric coupling}
the generated register \(K\), preserving the order of its twelve windows,
the phase origin and orientation. Its rational reading
\(\kappa=K_{\mathrm{int}}/(1000^{12}-1)\) is an encoding of that
register; it does not replace its types or the enriched history that produces it.
Hadamard inversion recovers its signed coordinate; the incidence
reading determines the Witt tetrad. The following constructions
make two further functions explicit: identifying operators on its cyclic
carrier and preserving information between two isometric realizations.
These are operations subsequent to the realization of the continuum, not rules
introducing conventional constants into APP--TRIT--TPK.

\subsection{Exact encoding and cyclic reference}
\label{iv:subsec:acoplamiento-codigo}

\begin{proposition}[Recovery of the record from its rational reading]
With base \(b=1000\), length twelve, origin, and orientation fixed, the map
\[
 (k_1,\ldots,k_{12})\longmapsto
 \kappa=\frac{\sum_{j=1}^{12}k_jb^{12-j}}{b^{12}-1},
 \qquad 0\leq k_j<b,
\]
is injective. From \(\kappa\), one recovers the integer
\(n=(b^{12}-1)\kappa\) and, by successive Euclidean divisions, the twelve
blocks, including leading zeros. An absolute error strictly
smaller than \(1/[2(b^{12}-1)]\) permits the same recovery by rounding.
\end{proposition}
\begin{proof}
The fixed-length integer expansion in base \(b\) is unique. Two distinct
integers are separated by at least one; their rational readings
are separated by at least \(1/(b^{12}-1)\). The strict bound places
the approximate datum within the rounding cell of a unique integer.
\end{proof}

Now consider \(K=(234,543,140,729,659,824,621,58,914,794,146,601)\),
the cyclic shift \(S\) of its positions, and
\(O_K=(K,SK,\ldots,S^{11}K)\) on \(\mathbb R^{12}\).
The shift of positions is identified neither with the entire enriched
evolution nor, solely by virtue of its periodicity, with an exceptional symmetry.

\begin{theorem}[Invertibility of the orbital reference]
\label{iv:thm:acoplamiento-orbita}
The matrix \(O_K\) is invertible, and
\[
 589609\|c\|^2\leq\|O_Kc\|^2\leq6263^2\|c\|^2.
\]
Its centred orbit has rank eleven.
\end{theorem}
\begin{proof}
The Gram matrix \(O_K^*O_K\) is circulant. Its first row, obtained from
the twelve inner products \(\langle K,S^jK\rangle\), is
\[
 \begin{aligned}
 (&4251257,2999323,3163385,3287920,3216553,3344743,\\
  &2950064,3344743,3216553,3287920,3163385,2999323).
 \end{aligned}
\]
Evaluation at the twelve characters of the cycle yields the
eigenvalues
\[
 \begin{gathered}
 6263^2,\quad697225,\quad
 589609\ (2),\quad1407529\ (2),\quad1053157\ (2),\\
 (1248025-345420\sqrt3)\ (2),\qquad
 (1248025+345420\sqrt3)\ (2),
 \end{gathered}
\]
where parentheses indicate multiplicities. The smallest value is
\(589609\): for the only comparison that is not immediate,
\(658416>345420\sqrt3\), as is seen by squaring both sides.
The triangle inequality for the Fourier transform and \(K_j\geq0\)
give the maximum \((\sum K_j)^2=6263^2\).
Positivity of every mode proves invertibility and the bounds.
Centring the record removes only the uniform mode.
\end{proof}

\begin{corollary}[Identification of an operator from its responses]
\label{iv:cor:acoplamiento-identificacion}
A linear operator \(T:\mathbb R^{12}\to\mathbb R^{12}\) is determined by
\(Y=(TK,TSK,\ldots,TS^{11}K)\), through \(T=YO_K^{-1}\).
If \(TS=ST\), the complete vector response \(y=TK\) suffices, since
\(T=O_yO_K^{-1}\). For a perturbation \(E\) of \(Y\),
\(\|\widehat T-T\|\leq\|E\|/\sqrt{589609}\).
\end{corollary}
\begin{proof}
We have \(Y=TO_K\). Under commutation, \(TS^jK=S^jy\), hence \(Y=O_y\).
The bound follows from
\(\|O_K^{-1}\|=1/\sqrt{589609}\).
The norms are the Euclidean norm and its induced operator norm.
\end{proof}

The response \(y\) contains twelve components. The statement does not
attribute the recovery of an arbitrary operator to a single scalar:
it specifies the observation and commutation condition that make identification possible.

\subsection{Conservation and reconstruction of two channels}
\label{iv:subsec:acoplamiento-bilateral}

Let \(H,H'\) be Hilbert spaces, and let
\(\mathcal B,\mathcal C:H\to H'\) be bounded linear isometries.
For \(a>0\), define
\[
 p=a^2+a+1,\qquad N=(2a+1)p=(a+1)^3+a^3,\qquad
 Q_-=a\mathcal B-\mathcal C,\quad Q_+=(a+1)\mathcal B+\mathcal C.
\]

\begin{theorem}[Bilateral balance and recovery]
\label{iv:thm:acoplamiento-balance}
The following hold
\[
 (a+1)Q_-^*Q_-+aQ_+^*Q_+=NI_H,\qquad
 W_av=\frac1{\sqrt N}
 \begin{pmatrix}\sqrt{a+1}\,Q_-v\\\sqrt a\,Q_+v\end{pmatrix},
 \qquad W_a^*W_a=I_H.
\]
The data \(x=Q_-v\), \(y=Q_+v\) determine
\[
 \mathcal Bv=\frac{x+y}{2a+1},\qquad
 \mathcal Cv=\frac{ay-(a+1)x}{2a+1},\qquad
 v=\mathcal B^*\frac{x+y}{2a+1}.
\]
\end{theorem}
\begin{proof}
Setting \(X=\mathcal B^*\mathcal C+\mathcal C^*\mathcal B\), one obtains
\[
 Q_-^*Q_-=(a^2+1)I-aX,\qquad
 Q_+^*Q_+=((a+1)^2+1)I+(a+1)X.
\]
The weights cancel \(X\), and the remaining coefficient is
\(2a^3+3a^2+3a+1=N\). Normalization gives the isometry.
Adding the two channels recovers \((2a+1)\mathcal Bv\);
the combination \(ay-(a+1)x\) recovers \((2a+1)\mathcal Cv\).
Finally, \(\mathcal B^*\mathcal B=I_H\).
\end{proof}

The balance also holds in infinite dimension: only bounded operators
are composed. The image of \(W_a\) is closed, and its orthogonal
projector is \(W_aW_a^*\). Thus a normalized pair has an
input precisely when it belongs to that image; its orthogonal
component measures the compatibility defect.

The factorization
\[
 W_a=R_aE_a,\qquad
 E_av=\frac1{\sqrt p}
 \begin{pmatrix}\sqrt{a(a+1)}\,\mathcal Bv\\\mathcal Cv\end{pmatrix},
 \quad
 R_a=\frac1{\sqrt{2a+1}}
 \begin{pmatrix}\sqrt a&-\sqrt{a+1}\\
 \sqrt{a+1}&\sqrt a\end{pmatrix}
\]
is verified by block multiplication; \(R_a\) acts on
\(H'\oplus H'\), with each scalar entry multiplying \(I_{H'}\).
In the nonadic partition
\(a=8\), the weights are \(72/73\) and \(1/73\); the factor
\(73=8\cdot9+1\) follows from the general law
\(a(a+1)+1\), rather than being an additional primitive datum. The corresponding balance is
\(9\|Q_-v\|^2+8\|Q_+v\|^2=17\cdot73\,\|v\|^2\).

\subsection{Transport of observables and incoming memory}
\label{iv:subsec:acoplamiento-lectores}

\begin{theorem}[Observables common to both channels]
\label{iv:thm:acoplamiento-lector}
For a bounded self-adjoint operator \(F\) on \(H'\),
\[
 W_a^*\operatorname{diag}(F,F)W_a
 =\frac{a(a+1)\mathcal B^*F\mathcal B+\mathcal C^*F\mathcal C}{p}.
\]
If \(\mathcal B\) is unitary and \(\mathcal C=\mathcal BR\), with
\(R\) unitary, and \(G=\mathcal B^*F\mathcal B\), the observable is
preserved if and only if \(GR=RG\).
\end{theorem}
\begin{proof}
Expanding \((a+1)Q_-^*FQ_-+aQ_+^*FQ_+\) cancels the cross terms.
The other two coefficients are
\(a(a+1)(2a+1)\) and \(2a+1\). Dividing by \(N\) gives the formula.
In the stated case, equality with \(G\) is equivalent to \(R^*GR=G\);
unitarity makes this equivalent to \(GR=RG\).
\end{proof}

The norm is obtained by taking \(F=I\). A general observable requires
the commutation condition: preserving the norm does not mean preserving
every observable without transformation.

\begin{theorem}[Unitary extension with memory]
\label{iv:thm:acoplamiento-memoria}
If \(H'=H\), \(\mathcal B=I\), and \(\mathcal C=U\) is unitary, set
\[
 A=\frac{\sqrt{a+1}}{\sqrt N}Q_-,\qquad
 D=\frac{\sqrt a}{\sqrt N}Q_+,\qquad
 \mathscr R_a=\begin{pmatrix}A&-D^*\\D&A^*\end{pmatrix}.
\]
Then \(\mathscr R_a\) is unitary on \(H\oplus H\), its first
column is \(W_a\), and it preserves \(\|v\|^2+\|m\|^2\) for every
input \((v,m)\). The inverse is
\[
 v=A^*z_1+D^*z_2,\qquad m=-Dz_1+Az_2.
\]
\end{theorem}
\begin{proof}
The operators \(A,D\) are normal and commute with each other and with their
adjoints, since they are polynomials in \(U,U^*\). The balance gives
\(A^*A+D^*D=I\), and normality gives \(AA^*+DD^*=I\).
The cross terms in both products
\(\mathscr R_a^*\mathscr R_a\) and \(\mathscr R_a\mathscr R_a^*\)
vanish; the diagonal blocks are the identity. Applying the adjoint to
\((z_1,z_2)\) yields the inverse.
\end{proof}

Recovery of the register, orbital identification and unitary
extension specify three distinct levels of information. The first
preserves the blocks with their chart; the second identifies operators under
the stated hypotheses; the third transports state and memory
jointly. This hierarchy allows the same generated structure to be used without
confusing a reduced reading with the state in its entirety.
